\documentclass[11pt]{article}
\usepackage[margin=1in]{geometry}
\usepackage[T1]{fontenc}
\usepackage{lmodern,amsmath,amssymb,amsthm,mathtools,booktabs}
\usepackage{xurl}
\usepackage[hidelinks]{hyperref}
\usepackage{microtype}
\setlength{\parindent}{0pt}
\setlength{\parskip}{0.55em}
\newtheorem{theorem}{Theorem}[section]
\newtheorem{lemma}[theorem]{Lemma}
\newtheorem{corollary}[theorem]{Corollary}
\newcommand{\K}{\mathcal K}
\newcommand{\R}{\mathbb R}
\newcommand{\norm}[1]{\left\lVert#1\right\rVert}
\newcommand{\diag}{\operatorname{diag}}
\newcommand{\rank}{\operatorname{rank}}
\newcommand{\range}{\operatorname{range}}
\newcommand{\Prob}{\mathbb P}
\begin{document}
\begin{center}
{\LARGE\bfseries RA-04: narrow-band stability\\[4pt] and exact-tail deflation}
\\[10pt]
{\large Sidney Holden}\\[3pt]
Center for Computational Biology, Flatiron Institute, Simons Foundation\\
{\small Flatiron Research Fellow, Biological Transport Networks}\\[3pt]
{\small Affiliation verified 13 September 2026 from the Simons Foundation staff profile.}
\end{center}
\begin{abstract}
This continuation establishes a quantitative good-start graph for sufficiently
narrow, nonzero-width leading spectral bands. Barycentric scaling eliminates
dependence on the overall leading spectral condition number. A separate generic
rank argument proves exact recovery when the spectrum below the target cutoff
has few distinct values, with the boundary multiplicity handled explicitly.
The unrestricted RA-04 bound is not established: the width restriction in the
first theorem is substantive, and the second theorem does not cover an arbitrary
diverse tail. The status is \textbf{PARTIAL}.
\end{abstract}
\tableofcontents

\section{Target and dependencies}
Let $M=AA^T$ have eigenvalues $\lambda_1\ge\cdots\ge\lambda_n\ge0$. Write
\[
 t=\lceil k/b\rceil,\qquad m=bt\le\rank(A),\qquad
 \Delta=\min_{1\le i\le m-b}\frac{\lambda_i-\lambda_{i+b}}{\lambda_i}>0.
\]
For $m=b$, the empty minimum is $\Delta=1$. The Gaussian starting block is
$G\in\R^{n\times b}$, and
\[
 \K_q(M,G)=\range[G,MG,\ldots,M^{q-1}G].
\]
RA-04 requests, for the prescribed projection-and-SVD-truncation output, all
three of its approximation guarantees with iteration order
\begin{equation}\label{target}
 q=O\!\left(\frac{t\log(2/\Delta)+\log(n/(\delta\varepsilon))}
 {\sqrt\varepsilon}\right).
\end{equation}
The previous report obtained the all-input bound with $t\log(2m/\Delta)$ in
place of $t\log(2/\Delta)$, together with several exact target regimes.\footnote{
\emph{RA04\_partial\_results.pdf}, supplied previous report; see also the
source inventory at the end of this report.}
The present work does not remove the additional $t\log m$ for arbitrary spectra.

The new graph theorem below is self-contained. Its implication for all three
RA-04 output guarantees uses the good-start convergence transfer of Chen et
al., Definition 3.1, Imported Theorem 3.2 and Observation 3.3 of arXiv v2, exactly as imported in the previous report.\footnote{T.~Chen,
E.~N.~Epperly, R.~A.~Meyer, C.~Musco, and A.~Rao, \emph{Does block size matter
in randomized block Krylov low-rank approximation?}, arXiv:2508.06486;
SODA 2026. \url{https://arxiv.org/abs/2508.06486}.}
In particular, a graph $U[I_m;F]$ at depth $q_0$ with
$\log(1+\norm{F}_F^2)=O(t\log(2/\Delta)+\log(n/\delta))$ and $q_0=O(t)$
provides the good start used in that transfer. No independent new proof of
the transfer theorem is claimed here. The fixed-spectrum corollary also uses
the earlier all-input bound; this additional dependency is stated where used.

\section{A Gaussian block estimate}
\begin{lemma}\label{gaussian}
Let $H\in\R^{a\times b}$ be standard Gaussian, $a\le b$. For $u\ge0$,
\[
 \Prob\{\sigma_{\min}(H)\le u\}\le16\sqrt a\,u.
\]
\end{lemma}
\begin{proof}
Let $v$ be a unit left singular vector for the smallest singular value, with
an independent fair sign. Orthogonal invariance makes $v$ uniform on the unit
sphere and independent of the singular values. For a fixed coordinate,
$\mathbb E v_i^2=1/a$ and $\mathbb E v_i^4=3/[a(a+2)]$. Thus
\[
 \Prob\{|v_i|\ge(2a)^{-1/2}\}\ge\frac1{12}
\]
by the Paley--Zygmund inequality. Repeated singular values have probability
zero and do not affect this assertion.

Let $D_i$ be the distance of row $i$ from the span of the other rows. The
Schur-complement identity gives
\[
 D_i^{-2}=((HH^T)^{-1})_{ii}\ge
 \frac{|v_i|^2}{\sigma_{\min}(H)^2}.
\]
Conditioned on the other rows, $D_i$ has a chi distribution with $b-a+1$
degrees of freedom. Consequently
$\Prob\{D_i\le z\}\le\sqrt{2/\pi}\,z$. Independence of $v$ from the
singular values now yields
\[
 \frac1{12}\Prob\{\sigma_{\min}(H)\le u\}
 \le\Prob\{D_i\le u\sqrt{2a}\}
 \le\frac2{\sqrt\pi}\sqrt a\,u.
\]
Since $24/\sqrt\pi<16$, the stated constant follows.
\end{proof}
For $r$ blocks with at most $b$ rows each, a union bound shows that every
least singular value is at least $\eta/(16r\sqrt b)$ except on an event of
probability at most $\eta$. Independence between this event and a subsequent
perturbation is not assumed or needed.

\section{Deterministic barycentric scaling}
Partition the first $m$ indices into $r\ge2$ consecutive nonempty groups
$J_s$ with $|J_s|\le b$. Let
\[
 c_1>\cdots>c_r>0,\qquad
 \beta=\min_{s<r}\left(1-\frac{c_{s+1}}{c_s}\right)>0,
 \qquad |\lambda_i-c_s|\le w c_s\quad(i\in J_s).
\]
Define the degree-$r-1$ polynomials and signed normalization constants
\[
 p_j(x)=\prod_{\ell\ne j}(x-c_\ell),\qquad
 d_s=\prod_{\ell\ne s}(c_s-c_\ell).
\]
\begin{lemma}\label{barycentric}
If $rw/\beta\le1/2$, then, for $i\in J_s$,
\begin{equation}\label{rowerror}
 \norm{\left(\frac{p_j(\lambda_i)}{d_s}\right)_{j=1}^r-e_s}
 _2\le\frac{2rw}{\beta}.
\end{equation}
Moreover,
\begin{equation}\label{tailratio}
 \frac{\displaystyle\max_{j,\,0\le x\le(1+w)c_r}|p_j(x)|}
 {\displaystyle\min_s|d_s|}
 \le\left(\frac{1+w}{\beta}\right)^{r-1}.
\end{equation}
Neither estimate contains $c_1/c_r$.
\end{lemma}
\begin{proof}
For any distinct $s,\ell$,
$|c_s-c_\ell|\ge\beta\max(c_s,c_\ell)$. The diagonal ratio is
\[
 \frac{p_s(\lambda_i)}{d_s}
 =\prod_{\ell\ne s}\left(1+
 \frac{\lambda_i-c_s}{c_s-c_\ell}\right).
\]
Each perturbation has absolute value at most $w/\beta$. Its difference from
one is at most $2(r-1)w/\beta$. For $j\ne s$, isolate the factor
$(\lambda_i-c_s)/(c_s-c_j)$; the other factors have product bounded by
$\exp((r-2)w/\beta)\le2$. The squared row error is therefore at most
$4[(r-1)^2+(r-1)](w/\beta)^2\le4r^2(w/\beta)^2$.

For the tail estimate, $|x-c_\ell|\le(1+w)c_\ell$, so
\[
 \max_j|p_j(x)|\le(1+w)^{r-1}\prod_{\ell<r}c_\ell.
\]
On the other hand, for every $s$,
\[
 |d_s|\ge\beta^{r-1}
 \left(\prod_{\ell<s}c_\ell\right)c_s^{r-s}
 \ge\beta^{r-1}\prod_{\ell<r}c_\ell.
\]
Division proves the claim. The cancellation is the reason to use this
normalization rather than perturb a raw monomial Vandermonde matrix.
\end{proof}

\section{A graph theorem for nonzero-width bands}
\begin{theorem}\label{graph}
Let $\eta=\delta/3$, where $0<\delta<1$. Under the preceding band assumptions,
suppose
\begin{equation}\label{width}
 w\le\frac{\beta\eta^{3/2}}{64r^2b\sqrt m}.
\end{equation}
With probability at least $1-\delta$, $\K_r(M,G)$ contains the columns of
$U[I_m;F]$, where $U$ is an eigenbasis of $M$ and
\begin{equation}\label{graphbound}
 \norm{F}_F\le32b r^{3/2}\sqrt{n-m}\,\eta^{-3/2}
 \left(\frac{1+w}{\beta}\right)^{r-1}.
\end{equation}
The empty tail when $n=m$ has norm zero.
\end{theorem}
\begin{proof}
In the eigenbasis write $U^TG=[H;T]$. The polynomial Krylov matrix with
basis $p_1,\ldots,p_r$ has head rows formed by concatenating
$p_j(\lambda_i)h_i^T$. Divide each row in group $s$ by $d_s$ and call the
result $C$. Its centered counterpart $C_0$ is block diagonal with blocks
$H_s\in\R^{|J_s|\times b}$. With $E=C-C_0$, Lemma~\ref{barycentric} gives
\[
 \norm E_2\le\norm E_F\le\frac{2rw}{\beta}\norm H_F.
\]
Intersect the events
\[
 \min_s\sigma_{\min}(H_s)\ge\frac\eta{16r\sqrt b},\qquad
 \norm H_F\le\sqrt{\frac{mb}\eta},\qquad
 \norm T_F\le\sqrt{\frac{(n-m)b}\eta}.
\]
Lemma~\ref{gaussian} and Markov's inequality bound total failure probability
by $3\eta=\delta$; omit the last event for an empty tail. On this intersection,
$R_0=C_0^\dagger$ satisfies $C_0R_0=I_m$ and
$\norm{R_0}_2\le16r\sqrt b/\eta$. The width bound implies
\[
 \norm{ER_0}_2\le
 \frac{32r^2b\sqrt m\,w}{\beta\eta^{3/2}}\le\frac12.
\]
Thus
\[
 R=R_0(I_m+ER_0)^{-1},\qquad CR=I_m,\qquad
 \norm R_2\le\frac{32r\sqrt b}\eta.
\]
This is a valid right inverse even for rectangular blocks.

Let $D$ have diagonal entry $d_s$ on $J_s$. The unscaled head is $DC$,
so $RD^{-1}$ is its right inverse. Multiplication of the full polynomial
Krylov matrix by this matrix produces head $I_m$ and tail $F$. The unscaled
tail has Frobenius norm at most
$\sqrt r\max_{x,j}|p_j(x)|\norm T_F$. Hence
\[
 \norm F_F\le\sqrt r\norm T_F\norm R_2
 \frac{\max_{x,j}|p_j(x)|}{\min_s|d_s|},
\]
which is \eqref{graphbound}. Each basis polynomial has degree at most $r-1$,
so the constructed space lies in $\K_r(M,G)$. Although $E$ and $C_0$ use the
same Gaussian rows, only deterministic inequalities on the stated event were
used. There is no independence error.
\end{proof}

\section{Consequences for the requested iteration order}
Take $r=t$ and $|J_s|=b$. Consecutive-band comparison gives
\[
 \beta-2w\le\Delta\le\beta+2w.
\]
In particular, $w\le\beta/4$ implies
$\beta/2\le\Delta\le3\beta/2$. Assumption~\eqref{width} is stronger than
this restriction. Since $b,r,m\le n$, Theorem~\ref{graph} yields
\[
 \log(1+\norm F_F^2)
 =O\bigl(t\log(2/\Delta)+\log(n/\delta)\bigr).
\]
The imported good-start transfer therefore gives \eqref{target} with all
three RA-04 guarantees. This proves the desired order for a nonzero-width
regime, not for arbitrary widths. Unequal bands also give this consequence
when $r=O(t)$ and their center gap is comparable to $\Delta$; the latter is
then an additional hypothesis.

\begin{corollary}\label{allfailure}
Suppose $m\ge2$, $t\ge2$, there are $t$ consecutive bands of size $b$, and
\begin{equation}\label{fixedwidth}
 w\le\frac{\beta}
 {64\,3^{3/2}t^2b\,m^{(3t+1)/2}}.
\end{equation}
Using the previous report's all-input bound, the RA-04 order
\eqref{target} holds for every $0<\delta<1$ and admissible $\varepsilon$.
\end{corollary}
\begin{proof}
When $\delta\ge m^{-t}$, condition~\eqref{fixedwidth} implies
\eqref{width}. When $\delta<m^{-t}$, the extra $t\log m$ in the previous
all-input bound is at most $\log(1/\delta)$ and is absorbed into the requested
separate logarithmic term. These two cases cover all failure probabilities.
\end{proof}
The previous all-input theorem is an explicit dependency of this corollary,
not a result re-proved in this continuation. No assertion is made that every
admissible spectrum satisfies \eqref{fixedwidth}.

\section{Generic block-Vandermonde rank}
\begin{theorem}\label{ranklemma}
Let $\rho$ rows be assigned real nodes, with at most $b$ rows at each node.
Let $H\in\R^{\rho\times b}$ have a joint density. If $qb\ge\rho$, then
\[
 [H,\Lambda H,\ldots,\Lambda^{q-1}H]
\]
has row rank $\rho$ almost surely.
\end{theorem}
\begin{proof}
Assign each equal-node group to distinct currently least-loaded colors among
$b$ colors. The maximum and minimum loads differ by at most one after each
group: increasing a subset of the minimum loads either preserves the two
adjacent load levels, or raises all minima and some maxima to the next two
levels. Consequently the final largest load is $\lceil\rho/b\rceil\le q$.

For a witness, set each row of $H$ equal to its color's coordinate vector.
After permutations, the displayed matrix splits into ordinary Vandermonde
blocks, one per color. Within a color the nodes are distinct and their number
is at most $q$. Each block has full row rank, exhibiting a nonzero
$\rho\times\rho$ minor. That minor is a nonzero polynomial in the entries
of $H$, whose zero set has Lebesgue measure zero.
\end{proof}
The coordinate-vector construction is a witness that a polynomial is not
identically zero, not a claim about the Gaussian realization itself.

\section{Exact recovery by distinct-tail deflation}
\begin{theorem}\label{exact}
Let $\theta=\lambda_k>0$. Let $a$ count eigenvalues strictly above $\theta$,
and let $d_\theta$ be the full multiplicity of $\theta$. Assume that every
higher eigenvalue has multiplicity at most $b$ and $k-a\le b$. Put
\[
 \rho=a+\min(d_\theta,b),\qquad
 L=\#\{\text{distinct eigenvalues strictly below }\theta\}.
\]
Zero is counted once in $L$ when present. Almost surely,
$\K_{L+\lceil\rho/b\rceil}(M,G)$ contains an optimal rank-$k$ left singular
subspace. The prescribed projection-and-truncation algorithm then returns
an exact best rank-$k$ approximation.
\end{theorem}
\begin{proof}
Let $f(x)=\prod_{\tau<\theta}(x-\tau)$, with one factor for each distinct
lower eigenvalue. The block $f(M)G$ has zero components below $\theta$ and
nonzero deterministic row scalings in every higher eigenspace and at $\theta$.
Each higher eigenspace has accessible dimension equal to its multiplicity
almost surely. At $\theta$ the accessible dimension is $\min(d_\theta,b)$.

Choose that many fixed coordinate rows from the $\theta$ eigenspace and
all rows from the higher eigenspaces. On these $\rho$ coordinates the Gaussian
block, after nonzero deterministic scalings, has a joint density. Each node
has at most $b$ chosen rows, so Theorem~\ref{ranklemma} gives full row rank
after $\lceil\rho/b\rceil$ blocks. If $d_\theta>b$, the remaining rows
are linear combinations of the chosen $b$ rows; since their eigenvalue is
identical, this relation holds for every power. Thus the full filtered space
has dimension exactly $\rho$ and spans all accessible directions.

It contains every higher eigenspace and at least $k-a$ directions at $\theta$,
which is an optimal rank-$k$ left singular subspace. Multiplication by $f$
adds $L$ to the maximum polynomial degree, giving the asserted number of
original Krylov blocks. Projection onto a space containing such a subspace,
followed by the prescribed SVD truncation, is an exact best rank-$k$
approximation.
\end{proof}
RA-04's positive clustered gap implies the multiplicity assumptions: more
than $b$ copies above the cutoff, or more than $b$ required copies at the
cutoff, would give $\lambda_i=\lambda_{i+b}$ within the first $m$ entries.
Also $a\le k-1\le m-1$, so $\rho\le m+b-1$. Hence $L+t+1$ blocks suffice.
This regime allows arbitrarily large lower-tail rank and multiplicities;
it restricts the number of distinct lower values, not their total count.

\subsection{A necessary boundary correction for exact recovery}
Take $b=2$, $k=m=4$, and eigenvalues $(5,4,3,2,2,2)$. Then $\Delta=2/5$,
$L=0$, and $\rho=3+2=5$. Three blocks suffice by Theorem~\ref{exact};
two generically do not. In $\K_2$, the block at eigenvalue $2$ is
$H_{\rm tail}(u+2v)$. Since $H_{\rm tail}$ has rank two almost surely,
a vector supported only above the cutoff must have $u+2v=0$. The intersection
with that higher eigenspace therefore has dimension at most two. Exact
rank-four approximation would require all three directions strictly above
$2$, which is impossible in this intersection. This is not a counterexample
to the big-$O$ assertion in RA-04.

\section{Validation and remaining obligations}
The exact checker uses rational arithmetic to test balanced-coloring rank
witnesses, polynomial normalization, right-inverse identities, nonzero-width
graph estimates, and the boundary example. The optional NumPy checker supplies
floating-point diagnostics. Execution results are recorded in the archive's
\texttt{results/} directory; no test count is assumed in this prose.

These finite checks do not replace the proofs or independently establish the
imported convergence theorem. They do not prove the unrestricted interpolation
estimate or certify a universal iteration constant. The width assumptions
remain essential to the presented argument. The prior general bound still
contains an extra $t\log m$ outside the proved regimes.

The sufficient square-depth interpolation estimate in the previous report is
still unproved here. Additional Krylov powers might bypass it; even a
counterexample to that stronger estimate would not by itself refute RA-04.
The appropriate status of this archive is \textbf{PARTIAL}, not SOLVED.

\section*{Sources}
\addcontentsline{toc}{section}{Sources}
\begin{enumerate}
\item Open Problems in Numerical Linear Algebra, RA-04 repository entry.
\url{https://github.com/ajt60gaibb/OpenProblemsInNLA/tree/main/randomized-and-low-rank-approximation/RA-04}.
\item \emph{RA04\_partial\_results.pdf}, previous report supplied in this
conversation. Used for the earlier all-input bound and the convergence-transfer
dependency. See the immutable prior submission and its independent audit at
\href{https://github.com/sidneyholden1/OpenProblemsInNLA/tree/e15ce4854d1a9c9fa457b93da78fca04febf200a/references/holden-ra04-2026-09-12}{prior submission (immutable revision e15ce48)}.
\item T.~Chen, E.~N.~Epperly, R.~A.~Meyer, C.~Musco, and A.~Rao,
\emph{Does block size matter in randomized block Krylov low-rank approximation?},
arXiv:2508.06486; SODA 2026. \url{https://arxiv.org/abs/2508.06486}.
\item R.~A.~Meyer, C.~Musco, and C.~Musco,
\emph{On the Unreasonable Effectiveness of Single Vector Krylov Methods for
Low-Rank Approximation}, arXiv:2305.02535.
\url{https://arxiv.org/abs/2305.02535}.
\item N.~Shao, \emph{A structural bound for cluster robustness of randomized
small-block Lanczos}.\newline arXiv:2507.10144.
\url{https://arxiv.org/abs/2507.10144}. Related work, not a proof of the
unresolved estimate.
\end{enumerate}
\end{document}
